% Figure: one time step as two sweeps (Chapter 5). Fractions, PLIC lines,
% donating regions, velocities and dilation volumes of the uniform fixture:
% generated/vof_sweeps.tex (vof.jl, sweeps_data: the independent sweep
% sequence, whose final fractions equal the code's). Labels hand-written.
\begin{figure}
  \centering
  \def\fxvnone{%
    \def\fxvcell##1##2##3##4##5##6##7{}\def\fxvliquid##1{}\def\fxvplic##1##2##3##4{}%
    \def\fxvsolid##1{}\def\fxvwall##1##2##3##4{}\def\fxvdr##1##2##3##4{}%
    \def\fxvvel##1##2##3##4{}\def\fxvdil##1##2##3{}}%
  \def\fxvpanelpicture#1{%
    \begin{tikzpicture}[x=0.92cm,y=0.92cm]
      \def\fxvpanel{#1}%
      \begin{scope}
        \clip (3,6) rectangle (7,9);
        \fxvnone
        \def\fxvliquid##1{\fill[tr liquid] ##1;}%
        % Generated by docs/techref/figures/generate.jl; do not edit.
\ifnum\fxvpanel=1
\fxvcell{3}{6}{4}{7}{0.19}{0}{0}
\fxvliquid{(3.654628,6) -- (4,6) -- (4,7) -- (3.973565,7) -- cycle}
\fxvplic{3.654628}{6}{3.973565}{7}
\fxvcell{4}{6}{5}{7}{1.00}{1}{0}
\fxvliquid{(4,6) -- (5,6) -- (5,7) -- (4,7) -- cycle}
\fxvcell{5}{6}{6}{7}{1.00}{1}{0}
\fxvliquid{(5,6) -- (6,6) -- (6,7) -- (5,7) -- cycle}
\fxvcell{6}{6}{7}{7}{1.00}{1}{0}
\fxvliquid{(6,6) -- (7,6) -- (7,7) -- (6,7) -- cycle}
\fxvcell{3}{7}{4}{8}{0.00}{0}{0}
\fxvliquid{(3.993273,7) -- (4,7) -- (4,7.012065) -- cycle}
\fxvplic{3.993273}{7}{4}{7.012065}
\fxvcell{4}{7}{5}{8}{0.63}{1}{0}
\fxvliquid{(4,7) -- (5,7) -- (5,7.992971) -- (4,7.263568) -- cycle}
\fxvplic{5}{7.992971}{4}{7.263568}
\fxvcell{5}{7}{6}{8}{1.00}{1}{0}
\fxvliquid{(5,7) -- (6,7) -- (6,8) -- (5,8) -- cycle}
\fxvcell{6}{7}{7}{8}{1.00}{1}{0}
\fxvliquid{(6,7) -- (7,7) -- (7,8) -- (6,8) -- cycle}
\fxvcell{3}{8}{4}{9}{0.00}{0}{0}
\fxvcell{4}{8}{5}{9}{0.00}{0}{0}
\fxvliquid{(4.915223,8) -- (5,8) -- (5,8.052949) -- cycle}
\fxvplic{4.915223}{8}{5}{8.052949}
\fxvcell{5}{8}{6}{9}{0.26}{0}{0}
\fxvliquid{(5,8) -- (6,8) -- (6,8.441461) -- (5,8.068701) -- cycle}
\fxvplic{6}{8.441461}{5}{8.068701}
\fxvcell{6}{8}{7}{9}{0.38}{0}{0}
\fxvliquid{(6,8) -- (7,8) -- (7,8.333053) -- (6,8.419014) -- cycle}
\fxvplic{7}{8.333053}{6}{8.419014}
\fxvdr{(3,6) -- (3.5,6) -- (3.5,7) -- (3,7) -- cycle}{(3,6) -- (3.5,6) -- (3.5,7) -- (3,7) -- cycle}{}{0.50}
\fxvdr{(3,7) -- (3.366025,7) -- (3.366025,8) -- (3,8) -- cycle}{(3,7) -- (3.366025,7) -- (3.366025,8) -- (3,8) -- cycle}{}{0.37}
\fxvdr{(3,8) -- (3.133975,8) -- (3.133975,9) -- (3,9) -- cycle}{(3,8) -- (3.133975,8) -- (3.133975,9) -- (3,9) -- cycle}{}{0.13}
\fxvdr{(4,6) -- (4.433013,6) -- (4.433013,7) -- (4,7) -- cycle}{(4,6) -- (4.433013,6) -- (4.433013,7) -- (4,7) -- cycle}{(4,6) -- (4.433013,6) -- (4.433013,7) -- (4,7) -- cycle}{0.43}
\fxvdr{(4,7) -- (4.316987,7) -- (4.316987,8) -- (4,8) -- cycle}{(4,7) -- (4.316987,7) -- (4.316987,8) -- (4,8) -- cycle}{(4,7) -- (4.316987,7) -- (4.316987,7.494779) -- (4,7.263568) -- cycle}{0.32}
\fxvdr{(4,8) -- (4.116025,8) -- (4.116025,9) -- (4,9) -- cycle}{(4,8) -- (4.116025,8) -- (4.116025,9) -- (4,9) -- cycle}{}{0.12}
\fxvdr{(5,6) -- (5.25,6) -- (5.25,7) -- (5,7) -- cycle}{(5,6) -- (5.25,6) -- (5.25,7) -- (5,7) -- cycle}{(5,6) -- (5.25,6) -- (5.25,7) -- (5,7) -- cycle}{0.25}
\fxvdr{(5,7) -- (5.183013,7) -- (5.183013,8) -- (5,8) -- cycle}{(5,7) -- (5.183013,7) -- (5.183013,8) -- (5,8) -- cycle}{(5,7) -- (5.183013,7) -- (5.183013,8) -- (5,8) -- cycle}{0.18}
\fxvdr{(5,8) -- (5.066987,8) -- (5.066987,9) -- (5,9) -- cycle}{(5,8) -- (5.066987,8) -- (5.066987,9) -- (5,9) -- cycle}{(5,8) -- (5.066987,8) -- (5.066987,8.093671) -- (5,8.068701) -- cycle}{0.07}
\fxvdr{(6.75,6) -- (7,6) -- (7,7) -- (6.75,7) -- cycle}{(6.75,6) -- (7,6) -- (7,7) -- (6.75,7) -- cycle}{(6.75,6) -- (7,6) -- (7,7) -- (6.75,7) -- cycle}{0.25}
\fxvdr{(6.816987,7) -- (7,7) -- (7,8) -- (6.816987,8) -- cycle}{(6.816987,7) -- (7,7) -- (7,8) -- (6.816987,8) -- cycle}{(6.816987,7) -- (7,7) -- (7,8) -- (6.816987,8) -- cycle}{0.18}
\fxvdr{(6.933013,8) -- (7,8) -- (7,9) -- (6.933013,9) -- cycle}{(6.933013,8) -- (7,8) -- (7,9) -- (6.933013,9) -- cycle}{(6.933013,8) -- (7,8) -- (7,8.333053) -- (6.933013,8.338811) -- cycle}{0.07}
\fxvvel{1}{3}{6.5}{-0.955}
\fxvvel{1}{3}{7.5}{-0.699}
\fxvvel{1}{3}{8.5}{-0.256}
\fxvvel{1}{4}{6.5}{-0.827}
\fxvvel{1}{4}{7.5}{-0.605}
\fxvvel{1}{4}{8.5}{-0.222}
\fxvvel{1}{5}{6.5}{-0.477}
\fxvvel{1}{5}{7.5}{-0.350}
\fxvvel{1}{5}{8.5}{-0.128}
\fxvvel{1}{7}{6.5}{0.477}
\fxvvel{1}{7}{7.5}{0.350}
\fxvvel{1}{7}{8.5}{0.128}
\fi
\ifnum\fxvpanel=2
\fxvcell{3}{6}{4}{7}{0.62}{0}{0}
\fxvliquid{(3.192609,6) -- (4,6) -- (4,7) -- (3.569558,7) -- cycle}
\fxvplic{3.192609}{6}{3.569558}{7}
\fxvcell{4}{6}{5}{7}{1.00}{1}{0}
\fxvliquid{(4,6) -- (5,6) -- (5,7) -- (4,7) -- cycle}
\fxvcell{5}{6}{6}{7}{1.00}{1}{0}
\fxvliquid{(5,6) -- (6,6) -- (6,7) -- (5,7) -- cycle}
\fxvcell{6}{6}{7}{7}{1.00}{1}{0}
\fxvliquid{(6,6) -- (7,6) -- (7,7) -- (6,7) -- cycle}
\fxvcell{3}{7}{4}{8}{0.12}{0}{0}
\fxvliquid{(3.597528,7) -- (4,7) -- (4,7.597477) -- cycle}
\fxvplic{3.597528}{7}{4}{7.597477}
\fxvcell{4}{7}{5}{8}{0.83}{1}{0}
\fxvliquid{(4,7) -- (5,7) -- (5,8) -- (4.565623,8) -- (4,7.381437) -- cycle}
\fxvplic{4.565623}{8}{4}{7.381437}
\fxvcell{5}{7}{6}{8}{1.00}{1}{0}
\fxvliquid{(5,7) -- (6,7) -- (6,8) -- (5,8) -- cycle}
\fxvcell{6}{7}{7}{8}{1.00}{1}{0}
\fxvliquid{(6,7) -- (7,7) -- (7,8) -- (6,8) -- cycle}
\fxvcell{3}{8}{4}{9}{0.00}{0}{0}
\fxvcell{4}{8}{5}{9}{0.01}{0}{0}
\fxvliquid{(4.808017,8) -- (5,8) -- (5,8.080037) -- cycle}
\fxvplic{4.808017}{8}{5}{8.080037}
\fxvcell{5}{8}{6}{9}{0.25}{0}{0}
\fxvliquid{(5,8) -- (6,8) -- (6,8.379838) -- (5,8.119447) -- cycle}
\fxvplic{6}{8.379838}{5}{8.119447}
\fxvcell{6}{8}{7}{9}{0.35}{0}{0}
\fxvliquid{(6,8) -- (7,8) -- (7,8.326195) -- (6,8.380865) -- cycle}
\fxvplic{7}{8.326195}{6}{8.380865}
\fxvdr{(3,7) -- (4,7) -- (4,7.066987) -- (3,7.066987) -- cycle}{(3,7) -- (4,7) -- (4,7.066987) -- (3,7.066987) -- cycle}{(3.597528,7) -- (4,7) -- (4,7.066987) -- (3.642652,7.066987) -- cycle}{0.07}
\fxvdr{(4,7) -- (5,7) -- (5,7.183013) -- (4,7.183013) -- cycle}{(4,7) -- (5,7) -- (5,7.183013) -- (4,7.183013) -- cycle}{(4,7) -- (5,7) -- (5,7.183013) -- (4,7.183013) -- cycle}{0.18}
\fxvdr{(5,7) -- (6,7) -- (6,7.25) -- (5,7.25) -- cycle}{(5,7) -- (6,7) -- (6,7.25) -- (5,7.25) -- cycle}{(5,7) -- (6,7) -- (6,7.25) -- (5,7.25) -- cycle}{0.25}
\fxvdr{(6,7) -- (7,7) -- (7,7.25) -- (6,7.25) -- cycle}{(6,7) -- (7,7) -- (7,7.25) -- (6,7.25) -- cycle}{(6,7) -- (7,7) -- (7,7.25) -- (6,7.25) -- cycle}{0.25}
\fxvdr{(3,8) -- (4,8) -- (4,8.116025) -- (3,8.116025) -- cycle}{(3,8) -- (4,8) -- (4,8.116025) -- (3,8.116025) -- cycle}{}{0.12}
\fxvdr{(4,8) -- (5,8) -- (5,8.316987) -- (4,8.316987) -- cycle}{(4,8) -- (5,8) -- (5,8.316987) -- (4,8.316987) -- cycle}{(4.808017,8) -- (5,8) -- (5,8.080037) -- cycle}{0.32}
\fxvdr{(5,8) -- (6,8) -- (6,8.433013) -- (5,8.433013) -- cycle}{(5,8) -- (6,8) -- (6,8.433013) -- (5,8.433013) -- cycle}{(5,8) -- (6,8) -- (6,8.379838) -- (5,8.119447) -- cycle}{0.43}
\fxvdr{(6,8) -- (7,8) -- (7,8.433013) -- (6,8.433013) -- cycle}{(6,8) -- (7,8) -- (7,8.433013) -- (6,8.433013) -- cycle}{(6,8) -- (7,8) -- (7,8.326195) -- (6,8.380865) -- cycle}{0.43}
\fxvvel{2}{3.5}{7}{-0.128}
\fxvvel{2}{4.5}{7}{-0.350}
\fxvvel{2}{5.5}{7}{-0.477}
\fxvvel{2}{6.5}{7}{-0.477}
\fxvvel{2}{3.5}{8}{-0.222}
\fxvvel{2}{4.5}{8}{-0.605}
\fxvvel{2}{5.5}{8}{-0.827}
\fxvvel{2}{6.5}{8}{-0.827}
\fxvvel{2}{3.5}{9}{-0.256}
\fxvvel{2}{4.5}{9}{-0.699}
\fxvvel{2}{5.5}{9}{-0.955}
\fxvvel{2}{6.5}{9}{-0.955}
\fxvdil{6}{6}{+0.250}
\fxvdil{5}{7}{+0.183}
\fxvdil{5}{6}{+0.250}
\fxvdil{6}{7}{+0.183}
\fxvdil{4}{7}{+0.134}
\fxvdil{4}{6}{+0.183}
\fi
\ifnum\fxvpanel=3
\fxvcell{3}{6}{4}{7}{0.64}{0}{0}
\fxvliquid{(3,6) -- (4,6) -- (4,7) -- (3.857511,7) -- (3,6.170542) -- cycle}
\fxvplic{3.857511}{7}{3}{6.170542}
\fxvcell{4}{6}{5}{7}{1.00}{1}{0}
\fxvliquid{(4,6) -- (5,6) -- (5,7) -- (4,7) -- cycle}
\fxvcell{5}{6}{6}{7}{1.00}{1}{0}
\fxvliquid{(5,6) -- (6,6) -- (6,7) -- (5,7) -- cycle}
\fxvcell{6}{6}{7}{7}{1.00}{1}{0}
\fxvliquid{(6,6) -- (7,6) -- (7,7) -- (6,7) -- cycle}
\fxvcell{3}{7}{4}{8}{0.09}{0}{0}
\fxvliquid{(3.505936,7) -- (4,7) -- (4,7.383694) -- cycle}
\fxvplic{3.505936}{7}{4}{7.383694}
\fxvcell{4}{7}{5}{8}{0.52}{1}{0}
\fxvliquid{(4,7) -- (5,7) -- (5,7.904064) -- (4,7.127455) -- cycle}
\fxvplic{5}{7.904064}{4}{7.127455}
\fxvcell{5}{7}{6}{8}{0.82}{1}{0}
\fxvliquid{(5,7) -- (6,7) -- (6,7.967065) -- (5,7.666194) -- cycle}
\fxvplic{6}{7.967065}{5}{7.666194}
\fxvcell{6}{7}{7}{8}{0.92}{1}{0}
\fxvliquid{(6,7) -- (7,7) -- (7,7.972461) -- (6,7.868574) -- cycle}
\fxvplic{7}{7.972461}{6}{7.868574}
\fxvcell{3}{8}{4}{9}{0.00}{0}{0}
\fxvcell{4}{8}{5}{9}{0.00}{0}{0}
\fxvcell{5}{8}{6}{9}{0.00}{0}{0}
\fxvcell{6}{8}{7}{9}{0.00}{0}{0}
\fxvdil{6}{6}{-0.250}
\fxvdil{5}{7}{-0.183}
\fxvdil{5}{6}{-0.250}
\fxvdil{6}{7}{-0.183}
\fxvdil{4}{7}{-0.134}
\fxvdil{4}{6}{-0.183}
\fi
%
        \fxvnone
        \def\fxvdr##1##2##3##4{\def\fxvarg{##3}\ifx\fxvarg\empty\else
          \fill[tr flux] ##3;\fi}%
        % Generated by docs/techref/figures/generate.jl; do not edit.
\ifnum\fxvpanel=1
\fxvcell{3}{6}{4}{7}{0.19}{0}{0}
\fxvliquid{(3.654628,6) -- (4,6) -- (4,7) -- (3.973565,7) -- cycle}
\fxvplic{3.654628}{6}{3.973565}{7}
\fxvcell{4}{6}{5}{7}{1.00}{1}{0}
\fxvliquid{(4,6) -- (5,6) -- (5,7) -- (4,7) -- cycle}
\fxvcell{5}{6}{6}{7}{1.00}{1}{0}
\fxvliquid{(5,6) -- (6,6) -- (6,7) -- (5,7) -- cycle}
\fxvcell{6}{6}{7}{7}{1.00}{1}{0}
\fxvliquid{(6,6) -- (7,6) -- (7,7) -- (6,7) -- cycle}
\fxvcell{3}{7}{4}{8}{0.00}{0}{0}
\fxvliquid{(3.993273,7) -- (4,7) -- (4,7.012065) -- cycle}
\fxvplic{3.993273}{7}{4}{7.012065}
\fxvcell{4}{7}{5}{8}{0.63}{1}{0}
\fxvliquid{(4,7) -- (5,7) -- (5,7.992971) -- (4,7.263568) -- cycle}
\fxvplic{5}{7.992971}{4}{7.263568}
\fxvcell{5}{7}{6}{8}{1.00}{1}{0}
\fxvliquid{(5,7) -- (6,7) -- (6,8) -- (5,8) -- cycle}
\fxvcell{6}{7}{7}{8}{1.00}{1}{0}
\fxvliquid{(6,7) -- (7,7) -- (7,8) -- (6,8) -- cycle}
\fxvcell{3}{8}{4}{9}{0.00}{0}{0}
\fxvcell{4}{8}{5}{9}{0.00}{0}{0}
\fxvliquid{(4.915223,8) -- (5,8) -- (5,8.052949) -- cycle}
\fxvplic{4.915223}{8}{5}{8.052949}
\fxvcell{5}{8}{6}{9}{0.26}{0}{0}
\fxvliquid{(5,8) -- (6,8) -- (6,8.441461) -- (5,8.068701) -- cycle}
\fxvplic{6}{8.441461}{5}{8.068701}
\fxvcell{6}{8}{7}{9}{0.38}{0}{0}
\fxvliquid{(6,8) -- (7,8) -- (7,8.333053) -- (6,8.419014) -- cycle}
\fxvplic{7}{8.333053}{6}{8.419014}
\fxvdr{(3,6) -- (3.5,6) -- (3.5,7) -- (3,7) -- cycle}{(3,6) -- (3.5,6) -- (3.5,7) -- (3,7) -- cycle}{}{0.50}
\fxvdr{(3,7) -- (3.366025,7) -- (3.366025,8) -- (3,8) -- cycle}{(3,7) -- (3.366025,7) -- (3.366025,8) -- (3,8) -- cycle}{}{0.37}
\fxvdr{(3,8) -- (3.133975,8) -- (3.133975,9) -- (3,9) -- cycle}{(3,8) -- (3.133975,8) -- (3.133975,9) -- (3,9) -- cycle}{}{0.13}
\fxvdr{(4,6) -- (4.433013,6) -- (4.433013,7) -- (4,7) -- cycle}{(4,6) -- (4.433013,6) -- (4.433013,7) -- (4,7) -- cycle}{(4,6) -- (4.433013,6) -- (4.433013,7) -- (4,7) -- cycle}{0.43}
\fxvdr{(4,7) -- (4.316987,7) -- (4.316987,8) -- (4,8) -- cycle}{(4,7) -- (4.316987,7) -- (4.316987,8) -- (4,8) -- cycle}{(4,7) -- (4.316987,7) -- (4.316987,7.494779) -- (4,7.263568) -- cycle}{0.32}
\fxvdr{(4,8) -- (4.116025,8) -- (4.116025,9) -- (4,9) -- cycle}{(4,8) -- (4.116025,8) -- (4.116025,9) -- (4,9) -- cycle}{}{0.12}
\fxvdr{(5,6) -- (5.25,6) -- (5.25,7) -- (5,7) -- cycle}{(5,6) -- (5.25,6) -- (5.25,7) -- (5,7) -- cycle}{(5,6) -- (5.25,6) -- (5.25,7) -- (5,7) -- cycle}{0.25}
\fxvdr{(5,7) -- (5.183013,7) -- (5.183013,8) -- (5,8) -- cycle}{(5,7) -- (5.183013,7) -- (5.183013,8) -- (5,8) -- cycle}{(5,7) -- (5.183013,7) -- (5.183013,8) -- (5,8) -- cycle}{0.18}
\fxvdr{(5,8) -- (5.066987,8) -- (5.066987,9) -- (5,9) -- cycle}{(5,8) -- (5.066987,8) -- (5.066987,9) -- (5,9) -- cycle}{(5,8) -- (5.066987,8) -- (5.066987,8.093671) -- (5,8.068701) -- cycle}{0.07}
\fxvdr{(6.75,6) -- (7,6) -- (7,7) -- (6.75,7) -- cycle}{(6.75,6) -- (7,6) -- (7,7) -- (6.75,7) -- cycle}{(6.75,6) -- (7,6) -- (7,7) -- (6.75,7) -- cycle}{0.25}
\fxvdr{(6.816987,7) -- (7,7) -- (7,8) -- (6.816987,8) -- cycle}{(6.816987,7) -- (7,7) -- (7,8) -- (6.816987,8) -- cycle}{(6.816987,7) -- (7,7) -- (7,8) -- (6.816987,8) -- cycle}{0.18}
\fxvdr{(6.933013,8) -- (7,8) -- (7,9) -- (6.933013,9) -- cycle}{(6.933013,8) -- (7,8) -- (7,9) -- (6.933013,9) -- cycle}{(6.933013,8) -- (7,8) -- (7,8.333053) -- (6.933013,8.338811) -- cycle}{0.07}
\fxvvel{1}{3}{6.5}{-0.955}
\fxvvel{1}{3}{7.5}{-0.699}
\fxvvel{1}{3}{8.5}{-0.256}
\fxvvel{1}{4}{6.5}{-0.827}
\fxvvel{1}{4}{7.5}{-0.605}
\fxvvel{1}{4}{8.5}{-0.222}
\fxvvel{1}{5}{6.5}{-0.477}
\fxvvel{1}{5}{7.5}{-0.350}
\fxvvel{1}{5}{8.5}{-0.128}
\fxvvel{1}{7}{6.5}{0.477}
\fxvvel{1}{7}{7.5}{0.350}
\fxvvel{1}{7}{8.5}{0.128}
\fi
\ifnum\fxvpanel=2
\fxvcell{3}{6}{4}{7}{0.62}{0}{0}
\fxvliquid{(3.192609,6) -- (4,6) -- (4,7) -- (3.569558,7) -- cycle}
\fxvplic{3.192609}{6}{3.569558}{7}
\fxvcell{4}{6}{5}{7}{1.00}{1}{0}
\fxvliquid{(4,6) -- (5,6) -- (5,7) -- (4,7) -- cycle}
\fxvcell{5}{6}{6}{7}{1.00}{1}{0}
\fxvliquid{(5,6) -- (6,6) -- (6,7) -- (5,7) -- cycle}
\fxvcell{6}{6}{7}{7}{1.00}{1}{0}
\fxvliquid{(6,6) -- (7,6) -- (7,7) -- (6,7) -- cycle}
\fxvcell{3}{7}{4}{8}{0.12}{0}{0}
\fxvliquid{(3.597528,7) -- (4,7) -- (4,7.597477) -- cycle}
\fxvplic{3.597528}{7}{4}{7.597477}
\fxvcell{4}{7}{5}{8}{0.83}{1}{0}
\fxvliquid{(4,7) -- (5,7) -- (5,8) -- (4.565623,8) -- (4,7.381437) -- cycle}
\fxvplic{4.565623}{8}{4}{7.381437}
\fxvcell{5}{7}{6}{8}{1.00}{1}{0}
\fxvliquid{(5,7) -- (6,7) -- (6,8) -- (5,8) -- cycle}
\fxvcell{6}{7}{7}{8}{1.00}{1}{0}
\fxvliquid{(6,7) -- (7,7) -- (7,8) -- (6,8) -- cycle}
\fxvcell{3}{8}{4}{9}{0.00}{0}{0}
\fxvcell{4}{8}{5}{9}{0.01}{0}{0}
\fxvliquid{(4.808017,8) -- (5,8) -- (5,8.080037) -- cycle}
\fxvplic{4.808017}{8}{5}{8.080037}
\fxvcell{5}{8}{6}{9}{0.25}{0}{0}
\fxvliquid{(5,8) -- (6,8) -- (6,8.379838) -- (5,8.119447) -- cycle}
\fxvplic{6}{8.379838}{5}{8.119447}
\fxvcell{6}{8}{7}{9}{0.35}{0}{0}
\fxvliquid{(6,8) -- (7,8) -- (7,8.326195) -- (6,8.380865) -- cycle}
\fxvplic{7}{8.326195}{6}{8.380865}
\fxvdr{(3,7) -- (4,7) -- (4,7.066987) -- (3,7.066987) -- cycle}{(3,7) -- (4,7) -- (4,7.066987) -- (3,7.066987) -- cycle}{(3.597528,7) -- (4,7) -- (4,7.066987) -- (3.642652,7.066987) -- cycle}{0.07}
\fxvdr{(4,7) -- (5,7) -- (5,7.183013) -- (4,7.183013) -- cycle}{(4,7) -- (5,7) -- (5,7.183013) -- (4,7.183013) -- cycle}{(4,7) -- (5,7) -- (5,7.183013) -- (4,7.183013) -- cycle}{0.18}
\fxvdr{(5,7) -- (6,7) -- (6,7.25) -- (5,7.25) -- cycle}{(5,7) -- (6,7) -- (6,7.25) -- (5,7.25) -- cycle}{(5,7) -- (6,7) -- (6,7.25) -- (5,7.25) -- cycle}{0.25}
\fxvdr{(6,7) -- (7,7) -- (7,7.25) -- (6,7.25) -- cycle}{(6,7) -- (7,7) -- (7,7.25) -- (6,7.25) -- cycle}{(6,7) -- (7,7) -- (7,7.25) -- (6,7.25) -- cycle}{0.25}
\fxvdr{(3,8) -- (4,8) -- (4,8.116025) -- (3,8.116025) -- cycle}{(3,8) -- (4,8) -- (4,8.116025) -- (3,8.116025) -- cycle}{}{0.12}
\fxvdr{(4,8) -- (5,8) -- (5,8.316987) -- (4,8.316987) -- cycle}{(4,8) -- (5,8) -- (5,8.316987) -- (4,8.316987) -- cycle}{(4.808017,8) -- (5,8) -- (5,8.080037) -- cycle}{0.32}
\fxvdr{(5,8) -- (6,8) -- (6,8.433013) -- (5,8.433013) -- cycle}{(5,8) -- (6,8) -- (6,8.433013) -- (5,8.433013) -- cycle}{(5,8) -- (6,8) -- (6,8.379838) -- (5,8.119447) -- cycle}{0.43}
\fxvdr{(6,8) -- (7,8) -- (7,8.433013) -- (6,8.433013) -- cycle}{(6,8) -- (7,8) -- (7,8.433013) -- (6,8.433013) -- cycle}{(6,8) -- (7,8) -- (7,8.326195) -- (6,8.380865) -- cycle}{0.43}
\fxvvel{2}{3.5}{7}{-0.128}
\fxvvel{2}{4.5}{7}{-0.350}
\fxvvel{2}{5.5}{7}{-0.477}
\fxvvel{2}{6.5}{7}{-0.477}
\fxvvel{2}{3.5}{8}{-0.222}
\fxvvel{2}{4.5}{8}{-0.605}
\fxvvel{2}{5.5}{8}{-0.827}
\fxvvel{2}{6.5}{8}{-0.827}
\fxvvel{2}{3.5}{9}{-0.256}
\fxvvel{2}{4.5}{9}{-0.699}
\fxvvel{2}{5.5}{9}{-0.955}
\fxvvel{2}{6.5}{9}{-0.955}
\fxvdil{6}{6}{+0.250}
\fxvdil{5}{7}{+0.183}
\fxvdil{5}{6}{+0.250}
\fxvdil{6}{7}{+0.183}
\fxvdil{4}{7}{+0.134}
\fxvdil{4}{6}{+0.183}
\fi
\ifnum\fxvpanel=3
\fxvcell{3}{6}{4}{7}{0.64}{0}{0}
\fxvliquid{(3,6) -- (4,6) -- (4,7) -- (3.857511,7) -- (3,6.170542) -- cycle}
\fxvplic{3.857511}{7}{3}{6.170542}
\fxvcell{4}{6}{5}{7}{1.00}{1}{0}
\fxvliquid{(4,6) -- (5,6) -- (5,7) -- (4,7) -- cycle}
\fxvcell{5}{6}{6}{7}{1.00}{1}{0}
\fxvliquid{(5,6) -- (6,6) -- (6,7) -- (5,7) -- cycle}
\fxvcell{6}{6}{7}{7}{1.00}{1}{0}
\fxvliquid{(6,6) -- (7,6) -- (7,7) -- (6,7) -- cycle}
\fxvcell{3}{7}{4}{8}{0.09}{0}{0}
\fxvliquid{(3.505936,7) -- (4,7) -- (4,7.383694) -- cycle}
\fxvplic{3.505936}{7}{4}{7.383694}
\fxvcell{4}{7}{5}{8}{0.52}{1}{0}
\fxvliquid{(4,7) -- (5,7) -- (5,7.904064) -- (4,7.127455) -- cycle}
\fxvplic{5}{7.904064}{4}{7.127455}
\fxvcell{5}{7}{6}{8}{0.82}{1}{0}
\fxvliquid{(5,7) -- (6,7) -- (6,7.967065) -- (5,7.666194) -- cycle}
\fxvplic{6}{7.967065}{5}{7.666194}
\fxvcell{6}{7}{7}{8}{0.92}{1}{0}
\fxvliquid{(6,7) -- (7,7) -- (7,7.972461) -- (6,7.868574) -- cycle}
\fxvplic{7}{7.972461}{6}{7.868574}
\fxvcell{3}{8}{4}{9}{0.00}{0}{0}
\fxvcell{4}{8}{5}{9}{0.00}{0}{0}
\fxvcell{5}{8}{6}{9}{0.00}{0}{0}
\fxvcell{6}{8}{7}{9}{0.00}{0}{0}
\fxvdil{6}{6}{-0.250}
\fxvdil{5}{7}{-0.183}
\fxvdil{5}{6}{-0.250}
\fxvdil{6}{7}{-0.183}
\fxvdil{4}{7}{-0.134}
\fxvdil{4}{6}{-0.183}
\fi
%
        \fxvnone
        \def\fxvcell##1##2##3##4##5##6##7{\draw[tr cell] (##1,##2) rectangle (##3,##4);}%
        \def\fxvplic##1##2##3##4{\draw[tr plic] (##1,##2) -- (##3,##4);}%
        % Generated by docs/techref/figures/generate.jl; do not edit.
\ifnum\fxvpanel=1
\fxvcell{3}{6}{4}{7}{0.19}{0}{0}
\fxvliquid{(3.654628,6) -- (4,6) -- (4,7) -- (3.973565,7) -- cycle}
\fxvplic{3.654628}{6}{3.973565}{7}
\fxvcell{4}{6}{5}{7}{1.00}{1}{0}
\fxvliquid{(4,6) -- (5,6) -- (5,7) -- (4,7) -- cycle}
\fxvcell{5}{6}{6}{7}{1.00}{1}{0}
\fxvliquid{(5,6) -- (6,6) -- (6,7) -- (5,7) -- cycle}
\fxvcell{6}{6}{7}{7}{1.00}{1}{0}
\fxvliquid{(6,6) -- (7,6) -- (7,7) -- (6,7) -- cycle}
\fxvcell{3}{7}{4}{8}{0.00}{0}{0}
\fxvliquid{(3.993273,7) -- (4,7) -- (4,7.012065) -- cycle}
\fxvplic{3.993273}{7}{4}{7.012065}
\fxvcell{4}{7}{5}{8}{0.63}{1}{0}
\fxvliquid{(4,7) -- (5,7) -- (5,7.992971) -- (4,7.263568) -- cycle}
\fxvplic{5}{7.992971}{4}{7.263568}
\fxvcell{5}{7}{6}{8}{1.00}{1}{0}
\fxvliquid{(5,7) -- (6,7) -- (6,8) -- (5,8) -- cycle}
\fxvcell{6}{7}{7}{8}{1.00}{1}{0}
\fxvliquid{(6,7) -- (7,7) -- (7,8) -- (6,8) -- cycle}
\fxvcell{3}{8}{4}{9}{0.00}{0}{0}
\fxvcell{4}{8}{5}{9}{0.00}{0}{0}
\fxvliquid{(4.915223,8) -- (5,8) -- (5,8.052949) -- cycle}
\fxvplic{4.915223}{8}{5}{8.052949}
\fxvcell{5}{8}{6}{9}{0.26}{0}{0}
\fxvliquid{(5,8) -- (6,8) -- (6,8.441461) -- (5,8.068701) -- cycle}
\fxvplic{6}{8.441461}{5}{8.068701}
\fxvcell{6}{8}{7}{9}{0.38}{0}{0}
\fxvliquid{(6,8) -- (7,8) -- (7,8.333053) -- (6,8.419014) -- cycle}
\fxvplic{7}{8.333053}{6}{8.419014}
\fxvdr{(3,6) -- (3.5,6) -- (3.5,7) -- (3,7) -- cycle}{(3,6) -- (3.5,6) -- (3.5,7) -- (3,7) -- cycle}{}{0.50}
\fxvdr{(3,7) -- (3.366025,7) -- (3.366025,8) -- (3,8) -- cycle}{(3,7) -- (3.366025,7) -- (3.366025,8) -- (3,8) -- cycle}{}{0.37}
\fxvdr{(3,8) -- (3.133975,8) -- (3.133975,9) -- (3,9) -- cycle}{(3,8) -- (3.133975,8) -- (3.133975,9) -- (3,9) -- cycle}{}{0.13}
\fxvdr{(4,6) -- (4.433013,6) -- (4.433013,7) -- (4,7) -- cycle}{(4,6) -- (4.433013,6) -- (4.433013,7) -- (4,7) -- cycle}{(4,6) -- (4.433013,6) -- (4.433013,7) -- (4,7) -- cycle}{0.43}
\fxvdr{(4,7) -- (4.316987,7) -- (4.316987,8) -- (4,8) -- cycle}{(4,7) -- (4.316987,7) -- (4.316987,8) -- (4,8) -- cycle}{(4,7) -- (4.316987,7) -- (4.316987,7.494779) -- (4,7.263568) -- cycle}{0.32}
\fxvdr{(4,8) -- (4.116025,8) -- (4.116025,9) -- (4,9) -- cycle}{(4,8) -- (4.116025,8) -- (4.116025,9) -- (4,9) -- cycle}{}{0.12}
\fxvdr{(5,6) -- (5.25,6) -- (5.25,7) -- (5,7) -- cycle}{(5,6) -- (5.25,6) -- (5.25,7) -- (5,7) -- cycle}{(5,6) -- (5.25,6) -- (5.25,7) -- (5,7) -- cycle}{0.25}
\fxvdr{(5,7) -- (5.183013,7) -- (5.183013,8) -- (5,8) -- cycle}{(5,7) -- (5.183013,7) -- (5.183013,8) -- (5,8) -- cycle}{(5,7) -- (5.183013,7) -- (5.183013,8) -- (5,8) -- cycle}{0.18}
\fxvdr{(5,8) -- (5.066987,8) -- (5.066987,9) -- (5,9) -- cycle}{(5,8) -- (5.066987,8) -- (5.066987,9) -- (5,9) -- cycle}{(5,8) -- (5.066987,8) -- (5.066987,8.093671) -- (5,8.068701) -- cycle}{0.07}
\fxvdr{(6.75,6) -- (7,6) -- (7,7) -- (6.75,7) -- cycle}{(6.75,6) -- (7,6) -- (7,7) -- (6.75,7) -- cycle}{(6.75,6) -- (7,6) -- (7,7) -- (6.75,7) -- cycle}{0.25}
\fxvdr{(6.816987,7) -- (7,7) -- (7,8) -- (6.816987,8) -- cycle}{(6.816987,7) -- (7,7) -- (7,8) -- (6.816987,8) -- cycle}{(6.816987,7) -- (7,7) -- (7,8) -- (6.816987,8) -- cycle}{0.18}
\fxvdr{(6.933013,8) -- (7,8) -- (7,9) -- (6.933013,9) -- cycle}{(6.933013,8) -- (7,8) -- (7,9) -- (6.933013,9) -- cycle}{(6.933013,8) -- (7,8) -- (7,8.333053) -- (6.933013,8.338811) -- cycle}{0.07}
\fxvvel{1}{3}{6.5}{-0.955}
\fxvvel{1}{3}{7.5}{-0.699}
\fxvvel{1}{3}{8.5}{-0.256}
\fxvvel{1}{4}{6.5}{-0.827}
\fxvvel{1}{4}{7.5}{-0.605}
\fxvvel{1}{4}{8.5}{-0.222}
\fxvvel{1}{5}{6.5}{-0.477}
\fxvvel{1}{5}{7.5}{-0.350}
\fxvvel{1}{5}{8.5}{-0.128}
\fxvvel{1}{7}{6.5}{0.477}
\fxvvel{1}{7}{7.5}{0.350}
\fxvvel{1}{7}{8.5}{0.128}
\fi
\ifnum\fxvpanel=2
\fxvcell{3}{6}{4}{7}{0.62}{0}{0}
\fxvliquid{(3.192609,6) -- (4,6) -- (4,7) -- (3.569558,7) -- cycle}
\fxvplic{3.192609}{6}{3.569558}{7}
\fxvcell{4}{6}{5}{7}{1.00}{1}{0}
\fxvliquid{(4,6) -- (5,6) -- (5,7) -- (4,7) -- cycle}
\fxvcell{5}{6}{6}{7}{1.00}{1}{0}
\fxvliquid{(5,6) -- (6,6) -- (6,7) -- (5,7) -- cycle}
\fxvcell{6}{6}{7}{7}{1.00}{1}{0}
\fxvliquid{(6,6) -- (7,6) -- (7,7) -- (6,7) -- cycle}
\fxvcell{3}{7}{4}{8}{0.12}{0}{0}
\fxvliquid{(3.597528,7) -- (4,7) -- (4,7.597477) -- cycle}
\fxvplic{3.597528}{7}{4}{7.597477}
\fxvcell{4}{7}{5}{8}{0.83}{1}{0}
\fxvliquid{(4,7) -- (5,7) -- (5,8) -- (4.565623,8) -- (4,7.381437) -- cycle}
\fxvplic{4.565623}{8}{4}{7.381437}
\fxvcell{5}{7}{6}{8}{1.00}{1}{0}
\fxvliquid{(5,7) -- (6,7) -- (6,8) -- (5,8) -- cycle}
\fxvcell{6}{7}{7}{8}{1.00}{1}{0}
\fxvliquid{(6,7) -- (7,7) -- (7,8) -- (6,8) -- cycle}
\fxvcell{3}{8}{4}{9}{0.00}{0}{0}
\fxvcell{4}{8}{5}{9}{0.01}{0}{0}
\fxvliquid{(4.808017,8) -- (5,8) -- (5,8.080037) -- cycle}
\fxvplic{4.808017}{8}{5}{8.080037}
\fxvcell{5}{8}{6}{9}{0.25}{0}{0}
\fxvliquid{(5,8) -- (6,8) -- (6,8.379838) -- (5,8.119447) -- cycle}
\fxvplic{6}{8.379838}{5}{8.119447}
\fxvcell{6}{8}{7}{9}{0.35}{0}{0}
\fxvliquid{(6,8) -- (7,8) -- (7,8.326195) -- (6,8.380865) -- cycle}
\fxvplic{7}{8.326195}{6}{8.380865}
\fxvdr{(3,7) -- (4,7) -- (4,7.066987) -- (3,7.066987) -- cycle}{(3,7) -- (4,7) -- (4,7.066987) -- (3,7.066987) -- cycle}{(3.597528,7) -- (4,7) -- (4,7.066987) -- (3.642652,7.066987) -- cycle}{0.07}
\fxvdr{(4,7) -- (5,7) -- (5,7.183013) -- (4,7.183013) -- cycle}{(4,7) -- (5,7) -- (5,7.183013) -- (4,7.183013) -- cycle}{(4,7) -- (5,7) -- (5,7.183013) -- (4,7.183013) -- cycle}{0.18}
\fxvdr{(5,7) -- (6,7) -- (6,7.25) -- (5,7.25) -- cycle}{(5,7) -- (6,7) -- (6,7.25) -- (5,7.25) -- cycle}{(5,7) -- (6,7) -- (6,7.25) -- (5,7.25) -- cycle}{0.25}
\fxvdr{(6,7) -- (7,7) -- (7,7.25) -- (6,7.25) -- cycle}{(6,7) -- (7,7) -- (7,7.25) -- (6,7.25) -- cycle}{(6,7) -- (7,7) -- (7,7.25) -- (6,7.25) -- cycle}{0.25}
\fxvdr{(3,8) -- (4,8) -- (4,8.116025) -- (3,8.116025) -- cycle}{(3,8) -- (4,8) -- (4,8.116025) -- (3,8.116025) -- cycle}{}{0.12}
\fxvdr{(4,8) -- (5,8) -- (5,8.316987) -- (4,8.316987) -- cycle}{(4,8) -- (5,8) -- (5,8.316987) -- (4,8.316987) -- cycle}{(4.808017,8) -- (5,8) -- (5,8.080037) -- cycle}{0.32}
\fxvdr{(5,8) -- (6,8) -- (6,8.433013) -- (5,8.433013) -- cycle}{(5,8) -- (6,8) -- (6,8.433013) -- (5,8.433013) -- cycle}{(5,8) -- (6,8) -- (6,8.379838) -- (5,8.119447) -- cycle}{0.43}
\fxvdr{(6,8) -- (7,8) -- (7,8.433013) -- (6,8.433013) -- cycle}{(6,8) -- (7,8) -- (7,8.433013) -- (6,8.433013) -- cycle}{(6,8) -- (7,8) -- (7,8.326195) -- (6,8.380865) -- cycle}{0.43}
\fxvvel{2}{3.5}{7}{-0.128}
\fxvvel{2}{4.5}{7}{-0.350}
\fxvvel{2}{5.5}{7}{-0.477}
\fxvvel{2}{6.5}{7}{-0.477}
\fxvvel{2}{3.5}{8}{-0.222}
\fxvvel{2}{4.5}{8}{-0.605}
\fxvvel{2}{5.5}{8}{-0.827}
\fxvvel{2}{6.5}{8}{-0.827}
\fxvvel{2}{3.5}{9}{-0.256}
\fxvvel{2}{4.5}{9}{-0.699}
\fxvvel{2}{5.5}{9}{-0.955}
\fxvvel{2}{6.5}{9}{-0.955}
\fxvdil{6}{6}{+0.250}
\fxvdil{5}{7}{+0.183}
\fxvdil{5}{6}{+0.250}
\fxvdil{6}{7}{+0.183}
\fxvdil{4}{7}{+0.134}
\fxvdil{4}{6}{+0.183}
\fi
\ifnum\fxvpanel=3
\fxvcell{3}{6}{4}{7}{0.64}{0}{0}
\fxvliquid{(3,6) -- (4,6) -- (4,7) -- (3.857511,7) -- (3,6.170542) -- cycle}
\fxvplic{3.857511}{7}{3}{6.170542}
\fxvcell{4}{6}{5}{7}{1.00}{1}{0}
\fxvliquid{(4,6) -- (5,6) -- (5,7) -- (4,7) -- cycle}
\fxvcell{5}{6}{6}{7}{1.00}{1}{0}
\fxvliquid{(5,6) -- (6,6) -- (6,7) -- (5,7) -- cycle}
\fxvcell{6}{6}{7}{7}{1.00}{1}{0}
\fxvliquid{(6,6) -- (7,6) -- (7,7) -- (6,7) -- cycle}
\fxvcell{3}{7}{4}{8}{0.09}{0}{0}
\fxvliquid{(3.505936,7) -- (4,7) -- (4,7.383694) -- cycle}
\fxvplic{3.505936}{7}{4}{7.383694}
\fxvcell{4}{7}{5}{8}{0.52}{1}{0}
\fxvliquid{(4,7) -- (5,7) -- (5,7.904064) -- (4,7.127455) -- cycle}
\fxvplic{5}{7.904064}{4}{7.127455}
\fxvcell{5}{7}{6}{8}{0.82}{1}{0}
\fxvliquid{(5,7) -- (6,7) -- (6,7.967065) -- (5,7.666194) -- cycle}
\fxvplic{6}{7.967065}{5}{7.666194}
\fxvcell{6}{7}{7}{8}{0.92}{1}{0}
\fxvliquid{(6,7) -- (7,7) -- (7,7.972461) -- (6,7.868574) -- cycle}
\fxvplic{7}{7.972461}{6}{7.868574}
\fxvcell{3}{8}{4}{9}{0.00}{0}{0}
\fxvcell{4}{8}{5}{9}{0.00}{0}{0}
\fxvcell{5}{8}{6}{9}{0.00}{0}{0}
\fxvcell{6}{8}{7}{9}{0.00}{0}{0}
\fxvdil{6}{6}{-0.250}
\fxvdil{5}{7}{-0.183}
\fxvdil{5}{6}{-0.250}
\fxvdil{6}{7}{-0.183}
\fxvdil{4}{7}{-0.134}
\fxvdil{4}{6}{-0.183}
\fi
%
        \fxvnone
        \def\fxvdr##1##2##3##4{\draw[tr dr] ##1;}%
        % Generated by docs/techref/figures/generate.jl; do not edit.
\ifnum\fxvpanel=1
\fxvcell{3}{6}{4}{7}{0.19}{0}{0}
\fxvliquid{(3.654628,6) -- (4,6) -- (4,7) -- (3.973565,7) -- cycle}
\fxvplic{3.654628}{6}{3.973565}{7}
\fxvcell{4}{6}{5}{7}{1.00}{1}{0}
\fxvliquid{(4,6) -- (5,6) -- (5,7) -- (4,7) -- cycle}
\fxvcell{5}{6}{6}{7}{1.00}{1}{0}
\fxvliquid{(5,6) -- (6,6) -- (6,7) -- (5,7) -- cycle}
\fxvcell{6}{6}{7}{7}{1.00}{1}{0}
\fxvliquid{(6,6) -- (7,6) -- (7,7) -- (6,7) -- cycle}
\fxvcell{3}{7}{4}{8}{0.00}{0}{0}
\fxvliquid{(3.993273,7) -- (4,7) -- (4,7.012065) -- cycle}
\fxvplic{3.993273}{7}{4}{7.012065}
\fxvcell{4}{7}{5}{8}{0.63}{1}{0}
\fxvliquid{(4,7) -- (5,7) -- (5,7.992971) -- (4,7.263568) -- cycle}
\fxvplic{5}{7.992971}{4}{7.263568}
\fxvcell{5}{7}{6}{8}{1.00}{1}{0}
\fxvliquid{(5,7) -- (6,7) -- (6,8) -- (5,8) -- cycle}
\fxvcell{6}{7}{7}{8}{1.00}{1}{0}
\fxvliquid{(6,7) -- (7,7) -- (7,8) -- (6,8) -- cycle}
\fxvcell{3}{8}{4}{9}{0.00}{0}{0}
\fxvcell{4}{8}{5}{9}{0.00}{0}{0}
\fxvliquid{(4.915223,8) -- (5,8) -- (5,8.052949) -- cycle}
\fxvplic{4.915223}{8}{5}{8.052949}
\fxvcell{5}{8}{6}{9}{0.26}{0}{0}
\fxvliquid{(5,8) -- (6,8) -- (6,8.441461) -- (5,8.068701) -- cycle}
\fxvplic{6}{8.441461}{5}{8.068701}
\fxvcell{6}{8}{7}{9}{0.38}{0}{0}
\fxvliquid{(6,8) -- (7,8) -- (7,8.333053) -- (6,8.419014) -- cycle}
\fxvplic{7}{8.333053}{6}{8.419014}
\fxvdr{(3,6) -- (3.5,6) -- (3.5,7) -- (3,7) -- cycle}{(3,6) -- (3.5,6) -- (3.5,7) -- (3,7) -- cycle}{}{0.50}
\fxvdr{(3,7) -- (3.366025,7) -- (3.366025,8) -- (3,8) -- cycle}{(3,7) -- (3.366025,7) -- (3.366025,8) -- (3,8) -- cycle}{}{0.37}
\fxvdr{(3,8) -- (3.133975,8) -- (3.133975,9) -- (3,9) -- cycle}{(3,8) -- (3.133975,8) -- (3.133975,9) -- (3,9) -- cycle}{}{0.13}
\fxvdr{(4,6) -- (4.433013,6) -- (4.433013,7) -- (4,7) -- cycle}{(4,6) -- (4.433013,6) -- (4.433013,7) -- (4,7) -- cycle}{(4,6) -- (4.433013,6) -- (4.433013,7) -- (4,7) -- cycle}{0.43}
\fxvdr{(4,7) -- (4.316987,7) -- (4.316987,8) -- (4,8) -- cycle}{(4,7) -- (4.316987,7) -- (4.316987,8) -- (4,8) -- cycle}{(4,7) -- (4.316987,7) -- (4.316987,7.494779) -- (4,7.263568) -- cycle}{0.32}
\fxvdr{(4,8) -- (4.116025,8) -- (4.116025,9) -- (4,9) -- cycle}{(4,8) -- (4.116025,8) -- (4.116025,9) -- (4,9) -- cycle}{}{0.12}
\fxvdr{(5,6) -- (5.25,6) -- (5.25,7) -- (5,7) -- cycle}{(5,6) -- (5.25,6) -- (5.25,7) -- (5,7) -- cycle}{(5,6) -- (5.25,6) -- (5.25,7) -- (5,7) -- cycle}{0.25}
\fxvdr{(5,7) -- (5.183013,7) -- (5.183013,8) -- (5,8) -- cycle}{(5,7) -- (5.183013,7) -- (5.183013,8) -- (5,8) -- cycle}{(5,7) -- (5.183013,7) -- (5.183013,8) -- (5,8) -- cycle}{0.18}
\fxvdr{(5,8) -- (5.066987,8) -- (5.066987,9) -- (5,9) -- cycle}{(5,8) -- (5.066987,8) -- (5.066987,9) -- (5,9) -- cycle}{(5,8) -- (5.066987,8) -- (5.066987,8.093671) -- (5,8.068701) -- cycle}{0.07}
\fxvdr{(6.75,6) -- (7,6) -- (7,7) -- (6.75,7) -- cycle}{(6.75,6) -- (7,6) -- (7,7) -- (6.75,7) -- cycle}{(6.75,6) -- (7,6) -- (7,7) -- (6.75,7) -- cycle}{0.25}
\fxvdr{(6.816987,7) -- (7,7) -- (7,8) -- (6.816987,8) -- cycle}{(6.816987,7) -- (7,7) -- (7,8) -- (6.816987,8) -- cycle}{(6.816987,7) -- (7,7) -- (7,8) -- (6.816987,8) -- cycle}{0.18}
\fxvdr{(6.933013,8) -- (7,8) -- (7,9) -- (6.933013,9) -- cycle}{(6.933013,8) -- (7,8) -- (7,9) -- (6.933013,9) -- cycle}{(6.933013,8) -- (7,8) -- (7,8.333053) -- (6.933013,8.338811) -- cycle}{0.07}
\fxvvel{1}{3}{6.5}{-0.955}
\fxvvel{1}{3}{7.5}{-0.699}
\fxvvel{1}{3}{8.5}{-0.256}
\fxvvel{1}{4}{6.5}{-0.827}
\fxvvel{1}{4}{7.5}{-0.605}
\fxvvel{1}{4}{8.5}{-0.222}
\fxvvel{1}{5}{6.5}{-0.477}
\fxvvel{1}{5}{7.5}{-0.350}
\fxvvel{1}{5}{8.5}{-0.128}
\fxvvel{1}{7}{6.5}{0.477}
\fxvvel{1}{7}{7.5}{0.350}
\fxvvel{1}{7}{8.5}{0.128}
\fi
\ifnum\fxvpanel=2
\fxvcell{3}{6}{4}{7}{0.62}{0}{0}
\fxvliquid{(3.192609,6) -- (4,6) -- (4,7) -- (3.569558,7) -- cycle}
\fxvplic{3.192609}{6}{3.569558}{7}
\fxvcell{4}{6}{5}{7}{1.00}{1}{0}
\fxvliquid{(4,6) -- (5,6) -- (5,7) -- (4,7) -- cycle}
\fxvcell{5}{6}{6}{7}{1.00}{1}{0}
\fxvliquid{(5,6) -- (6,6) -- (6,7) -- (5,7) -- cycle}
\fxvcell{6}{6}{7}{7}{1.00}{1}{0}
\fxvliquid{(6,6) -- (7,6) -- (7,7) -- (6,7) -- cycle}
\fxvcell{3}{7}{4}{8}{0.12}{0}{0}
\fxvliquid{(3.597528,7) -- (4,7) -- (4,7.597477) -- cycle}
\fxvplic{3.597528}{7}{4}{7.597477}
\fxvcell{4}{7}{5}{8}{0.83}{1}{0}
\fxvliquid{(4,7) -- (5,7) -- (5,8) -- (4.565623,8) -- (4,7.381437) -- cycle}
\fxvplic{4.565623}{8}{4}{7.381437}
\fxvcell{5}{7}{6}{8}{1.00}{1}{0}
\fxvliquid{(5,7) -- (6,7) -- (6,8) -- (5,8) -- cycle}
\fxvcell{6}{7}{7}{8}{1.00}{1}{0}
\fxvliquid{(6,7) -- (7,7) -- (7,8) -- (6,8) -- cycle}
\fxvcell{3}{8}{4}{9}{0.00}{0}{0}
\fxvcell{4}{8}{5}{9}{0.01}{0}{0}
\fxvliquid{(4.808017,8) -- (5,8) -- (5,8.080037) -- cycle}
\fxvplic{4.808017}{8}{5}{8.080037}
\fxvcell{5}{8}{6}{9}{0.25}{0}{0}
\fxvliquid{(5,8) -- (6,8) -- (6,8.379838) -- (5,8.119447) -- cycle}
\fxvplic{6}{8.379838}{5}{8.119447}
\fxvcell{6}{8}{7}{9}{0.35}{0}{0}
\fxvliquid{(6,8) -- (7,8) -- (7,8.326195) -- (6,8.380865) -- cycle}
\fxvplic{7}{8.326195}{6}{8.380865}
\fxvdr{(3,7) -- (4,7) -- (4,7.066987) -- (3,7.066987) -- cycle}{(3,7) -- (4,7) -- (4,7.066987) -- (3,7.066987) -- cycle}{(3.597528,7) -- (4,7) -- (4,7.066987) -- (3.642652,7.066987) -- cycle}{0.07}
\fxvdr{(4,7) -- (5,7) -- (5,7.183013) -- (4,7.183013) -- cycle}{(4,7) -- (5,7) -- (5,7.183013) -- (4,7.183013) -- cycle}{(4,7) -- (5,7) -- (5,7.183013) -- (4,7.183013) -- cycle}{0.18}
\fxvdr{(5,7) -- (6,7) -- (6,7.25) -- (5,7.25) -- cycle}{(5,7) -- (6,7) -- (6,7.25) -- (5,7.25) -- cycle}{(5,7) -- (6,7) -- (6,7.25) -- (5,7.25) -- cycle}{0.25}
\fxvdr{(6,7) -- (7,7) -- (7,7.25) -- (6,7.25) -- cycle}{(6,7) -- (7,7) -- (7,7.25) -- (6,7.25) -- cycle}{(6,7) -- (7,7) -- (7,7.25) -- (6,7.25) -- cycle}{0.25}
\fxvdr{(3,8) -- (4,8) -- (4,8.116025) -- (3,8.116025) -- cycle}{(3,8) -- (4,8) -- (4,8.116025) -- (3,8.116025) -- cycle}{}{0.12}
\fxvdr{(4,8) -- (5,8) -- (5,8.316987) -- (4,8.316987) -- cycle}{(4,8) -- (5,8) -- (5,8.316987) -- (4,8.316987) -- cycle}{(4.808017,8) -- (5,8) -- (5,8.080037) -- cycle}{0.32}
\fxvdr{(5,8) -- (6,8) -- (6,8.433013) -- (5,8.433013) -- cycle}{(5,8) -- (6,8) -- (6,8.433013) -- (5,8.433013) -- cycle}{(5,8) -- (6,8) -- (6,8.379838) -- (5,8.119447) -- cycle}{0.43}
\fxvdr{(6,8) -- (7,8) -- (7,8.433013) -- (6,8.433013) -- cycle}{(6,8) -- (7,8) -- (7,8.433013) -- (6,8.433013) -- cycle}{(6,8) -- (7,8) -- (7,8.326195) -- (6,8.380865) -- cycle}{0.43}
\fxvvel{2}{3.5}{7}{-0.128}
\fxvvel{2}{4.5}{7}{-0.350}
\fxvvel{2}{5.5}{7}{-0.477}
\fxvvel{2}{6.5}{7}{-0.477}
\fxvvel{2}{3.5}{8}{-0.222}
\fxvvel{2}{4.5}{8}{-0.605}
\fxvvel{2}{5.5}{8}{-0.827}
\fxvvel{2}{6.5}{8}{-0.827}
\fxvvel{2}{3.5}{9}{-0.256}
\fxvvel{2}{4.5}{9}{-0.699}
\fxvvel{2}{5.5}{9}{-0.955}
\fxvvel{2}{6.5}{9}{-0.955}
\fxvdil{6}{6}{+0.250}
\fxvdil{5}{7}{+0.183}
\fxvdil{5}{6}{+0.250}
\fxvdil{6}{7}{+0.183}
\fxvdil{4}{7}{+0.134}
\fxvdil{4}{6}{+0.183}
\fi
\ifnum\fxvpanel=3
\fxvcell{3}{6}{4}{7}{0.64}{0}{0}
\fxvliquid{(3,6) -- (4,6) -- (4,7) -- (3.857511,7) -- (3,6.170542) -- cycle}
\fxvplic{3.857511}{7}{3}{6.170542}
\fxvcell{4}{6}{5}{7}{1.00}{1}{0}
\fxvliquid{(4,6) -- (5,6) -- (5,7) -- (4,7) -- cycle}
\fxvcell{5}{6}{6}{7}{1.00}{1}{0}
\fxvliquid{(5,6) -- (6,6) -- (6,7) -- (5,7) -- cycle}
\fxvcell{6}{6}{7}{7}{1.00}{1}{0}
\fxvliquid{(6,6) -- (7,6) -- (7,7) -- (6,7) -- cycle}
\fxvcell{3}{7}{4}{8}{0.09}{0}{0}
\fxvliquid{(3.505936,7) -- (4,7) -- (4,7.383694) -- cycle}
\fxvplic{3.505936}{7}{4}{7.383694}
\fxvcell{4}{7}{5}{8}{0.52}{1}{0}
\fxvliquid{(4,7) -- (5,7) -- (5,7.904064) -- (4,7.127455) -- cycle}
\fxvplic{5}{7.904064}{4}{7.127455}
\fxvcell{5}{7}{6}{8}{0.82}{1}{0}
\fxvliquid{(5,7) -- (6,7) -- (6,7.967065) -- (5,7.666194) -- cycle}
\fxvplic{6}{7.967065}{5}{7.666194}
\fxvcell{6}{7}{7}{8}{0.92}{1}{0}
\fxvliquid{(6,7) -- (7,7) -- (7,7.972461) -- (6,7.868574) -- cycle}
\fxvplic{7}{7.972461}{6}{7.868574}
\fxvcell{3}{8}{4}{9}{0.00}{0}{0}
\fxvcell{4}{8}{5}{9}{0.00}{0}{0}
\fxvcell{5}{8}{6}{9}{0.00}{0}{0}
\fxvcell{6}{8}{7}{9}{0.00}{0}{0}
\fxvdil{6}{6}{-0.250}
\fxvdil{5}{7}{-0.183}
\fxvdil{5}{6}{-0.250}
\fxvdil{6}{7}{-0.183}
\fxvdil{4}{7}{-0.134}
\fxvdil{4}{6}{-0.183}
\fi
%
      \end{scope}
      \fxvnone
      \def\fxvcell##1##2##3##4##5##6##7{%
        \node[tr value,font=\scriptsize] at ($(##1,##2)!0.5!(##3,##4)$) {##5};
        \ifnum##6=1 \node[tr chi] at ($(##1,##4)+(0.1,-0.1)$) {};\fi}%
      \def\fxvdil##1##2##3{\node[tr tiny,anchor=south west,text=trpurple] at (##1,##2) {$##3$};}%
      \def\fxvvel##1##2##3##4{\ifnum##1=1
          \draw[tr thin vector,trred] (##2,##3) -- ++(0.28*##4,0);
        \else
          \draw[tr thin vector,trred] (##2,##3) -- ++(0,0.28*##4);
        \fi}%
      % Generated by docs/techref/figures/generate.jl; do not edit.
\ifnum\fxvpanel=1
\fxvcell{3}{6}{4}{7}{0.19}{0}{0}
\fxvliquid{(3.654628,6) -- (4,6) -- (4,7) -- (3.973565,7) -- cycle}
\fxvplic{3.654628}{6}{3.973565}{7}
\fxvcell{4}{6}{5}{7}{1.00}{1}{0}
\fxvliquid{(4,6) -- (5,6) -- (5,7) -- (4,7) -- cycle}
\fxvcell{5}{6}{6}{7}{1.00}{1}{0}
\fxvliquid{(5,6) -- (6,6) -- (6,7) -- (5,7) -- cycle}
\fxvcell{6}{6}{7}{7}{1.00}{1}{0}
\fxvliquid{(6,6) -- (7,6) -- (7,7) -- (6,7) -- cycle}
\fxvcell{3}{7}{4}{8}{0.00}{0}{0}
\fxvliquid{(3.993273,7) -- (4,7) -- (4,7.012065) -- cycle}
\fxvplic{3.993273}{7}{4}{7.012065}
\fxvcell{4}{7}{5}{8}{0.63}{1}{0}
\fxvliquid{(4,7) -- (5,7) -- (5,7.992971) -- (4,7.263568) -- cycle}
\fxvplic{5}{7.992971}{4}{7.263568}
\fxvcell{5}{7}{6}{8}{1.00}{1}{0}
\fxvliquid{(5,7) -- (6,7) -- (6,8) -- (5,8) -- cycle}
\fxvcell{6}{7}{7}{8}{1.00}{1}{0}
\fxvliquid{(6,7) -- (7,7) -- (7,8) -- (6,8) -- cycle}
\fxvcell{3}{8}{4}{9}{0.00}{0}{0}
\fxvcell{4}{8}{5}{9}{0.00}{0}{0}
\fxvliquid{(4.915223,8) -- (5,8) -- (5,8.052949) -- cycle}
\fxvplic{4.915223}{8}{5}{8.052949}
\fxvcell{5}{8}{6}{9}{0.26}{0}{0}
\fxvliquid{(5,8) -- (6,8) -- (6,8.441461) -- (5,8.068701) -- cycle}
\fxvplic{6}{8.441461}{5}{8.068701}
\fxvcell{6}{8}{7}{9}{0.38}{0}{0}
\fxvliquid{(6,8) -- (7,8) -- (7,8.333053) -- (6,8.419014) -- cycle}
\fxvplic{7}{8.333053}{6}{8.419014}
\fxvdr{(3,6) -- (3.5,6) -- (3.5,7) -- (3,7) -- cycle}{(3,6) -- (3.5,6) -- (3.5,7) -- (3,7) -- cycle}{}{0.50}
\fxvdr{(3,7) -- (3.366025,7) -- (3.366025,8) -- (3,8) -- cycle}{(3,7) -- (3.366025,7) -- (3.366025,8) -- (3,8) -- cycle}{}{0.37}
\fxvdr{(3,8) -- (3.133975,8) -- (3.133975,9) -- (3,9) -- cycle}{(3,8) -- (3.133975,8) -- (3.133975,9) -- (3,9) -- cycle}{}{0.13}
\fxvdr{(4,6) -- (4.433013,6) -- (4.433013,7) -- (4,7) -- cycle}{(4,6) -- (4.433013,6) -- (4.433013,7) -- (4,7) -- cycle}{(4,6) -- (4.433013,6) -- (4.433013,7) -- (4,7) -- cycle}{0.43}
\fxvdr{(4,7) -- (4.316987,7) -- (4.316987,8) -- (4,8) -- cycle}{(4,7) -- (4.316987,7) -- (4.316987,8) -- (4,8) -- cycle}{(4,7) -- (4.316987,7) -- (4.316987,7.494779) -- (4,7.263568) -- cycle}{0.32}
\fxvdr{(4,8) -- (4.116025,8) -- (4.116025,9) -- (4,9) -- cycle}{(4,8) -- (4.116025,8) -- (4.116025,9) -- (4,9) -- cycle}{}{0.12}
\fxvdr{(5,6) -- (5.25,6) -- (5.25,7) -- (5,7) -- cycle}{(5,6) -- (5.25,6) -- (5.25,7) -- (5,7) -- cycle}{(5,6) -- (5.25,6) -- (5.25,7) -- (5,7) -- cycle}{0.25}
\fxvdr{(5,7) -- (5.183013,7) -- (5.183013,8) -- (5,8) -- cycle}{(5,7) -- (5.183013,7) -- (5.183013,8) -- (5,8) -- cycle}{(5,7) -- (5.183013,7) -- (5.183013,8) -- (5,8) -- cycle}{0.18}
\fxvdr{(5,8) -- (5.066987,8) -- (5.066987,9) -- (5,9) -- cycle}{(5,8) -- (5.066987,8) -- (5.066987,9) -- (5,9) -- cycle}{(5,8) -- (5.066987,8) -- (5.066987,8.093671) -- (5,8.068701) -- cycle}{0.07}
\fxvdr{(6.75,6) -- (7,6) -- (7,7) -- (6.75,7) -- cycle}{(6.75,6) -- (7,6) -- (7,7) -- (6.75,7) -- cycle}{(6.75,6) -- (7,6) -- (7,7) -- (6.75,7) -- cycle}{0.25}
\fxvdr{(6.816987,7) -- (7,7) -- (7,8) -- (6.816987,8) -- cycle}{(6.816987,7) -- (7,7) -- (7,8) -- (6.816987,8) -- cycle}{(6.816987,7) -- (7,7) -- (7,8) -- (6.816987,8) -- cycle}{0.18}
\fxvdr{(6.933013,8) -- (7,8) -- (7,9) -- (6.933013,9) -- cycle}{(6.933013,8) -- (7,8) -- (7,9) -- (6.933013,9) -- cycle}{(6.933013,8) -- (7,8) -- (7,8.333053) -- (6.933013,8.338811) -- cycle}{0.07}
\fxvvel{1}{3}{6.5}{-0.955}
\fxvvel{1}{3}{7.5}{-0.699}
\fxvvel{1}{3}{8.5}{-0.256}
\fxvvel{1}{4}{6.5}{-0.827}
\fxvvel{1}{4}{7.5}{-0.605}
\fxvvel{1}{4}{8.5}{-0.222}
\fxvvel{1}{5}{6.5}{-0.477}
\fxvvel{1}{5}{7.5}{-0.350}
\fxvvel{1}{5}{8.5}{-0.128}
\fxvvel{1}{7}{6.5}{0.477}
\fxvvel{1}{7}{7.5}{0.350}
\fxvvel{1}{7}{8.5}{0.128}
\fi
\ifnum\fxvpanel=2
\fxvcell{3}{6}{4}{7}{0.62}{0}{0}
\fxvliquid{(3.192609,6) -- (4,6) -- (4,7) -- (3.569558,7) -- cycle}
\fxvplic{3.192609}{6}{3.569558}{7}
\fxvcell{4}{6}{5}{7}{1.00}{1}{0}
\fxvliquid{(4,6) -- (5,6) -- (5,7) -- (4,7) -- cycle}
\fxvcell{5}{6}{6}{7}{1.00}{1}{0}
\fxvliquid{(5,6) -- (6,6) -- (6,7) -- (5,7) -- cycle}
\fxvcell{6}{6}{7}{7}{1.00}{1}{0}
\fxvliquid{(6,6) -- (7,6) -- (7,7) -- (6,7) -- cycle}
\fxvcell{3}{7}{4}{8}{0.12}{0}{0}
\fxvliquid{(3.597528,7) -- (4,7) -- (4,7.597477) -- cycle}
\fxvplic{3.597528}{7}{4}{7.597477}
\fxvcell{4}{7}{5}{8}{0.83}{1}{0}
\fxvliquid{(4,7) -- (5,7) -- (5,8) -- (4.565623,8) -- (4,7.381437) -- cycle}
\fxvplic{4.565623}{8}{4}{7.381437}
\fxvcell{5}{7}{6}{8}{1.00}{1}{0}
\fxvliquid{(5,7) -- (6,7) -- (6,8) -- (5,8) -- cycle}
\fxvcell{6}{7}{7}{8}{1.00}{1}{0}
\fxvliquid{(6,7) -- (7,7) -- (7,8) -- (6,8) -- cycle}
\fxvcell{3}{8}{4}{9}{0.00}{0}{0}
\fxvcell{4}{8}{5}{9}{0.01}{0}{0}
\fxvliquid{(4.808017,8) -- (5,8) -- (5,8.080037) -- cycle}
\fxvplic{4.808017}{8}{5}{8.080037}
\fxvcell{5}{8}{6}{9}{0.25}{0}{0}
\fxvliquid{(5,8) -- (6,8) -- (6,8.379838) -- (5,8.119447) -- cycle}
\fxvplic{6}{8.379838}{5}{8.119447}
\fxvcell{6}{8}{7}{9}{0.35}{0}{0}
\fxvliquid{(6,8) -- (7,8) -- (7,8.326195) -- (6,8.380865) -- cycle}
\fxvplic{7}{8.326195}{6}{8.380865}
\fxvdr{(3,7) -- (4,7) -- (4,7.066987) -- (3,7.066987) -- cycle}{(3,7) -- (4,7) -- (4,7.066987) -- (3,7.066987) -- cycle}{(3.597528,7) -- (4,7) -- (4,7.066987) -- (3.642652,7.066987) -- cycle}{0.07}
\fxvdr{(4,7) -- (5,7) -- (5,7.183013) -- (4,7.183013) -- cycle}{(4,7) -- (5,7) -- (5,7.183013) -- (4,7.183013) -- cycle}{(4,7) -- (5,7) -- (5,7.183013) -- (4,7.183013) -- cycle}{0.18}
\fxvdr{(5,7) -- (6,7) -- (6,7.25) -- (5,7.25) -- cycle}{(5,7) -- (6,7) -- (6,7.25) -- (5,7.25) -- cycle}{(5,7) -- (6,7) -- (6,7.25) -- (5,7.25) -- cycle}{0.25}
\fxvdr{(6,7) -- (7,7) -- (7,7.25) -- (6,7.25) -- cycle}{(6,7) -- (7,7) -- (7,7.25) -- (6,7.25) -- cycle}{(6,7) -- (7,7) -- (7,7.25) -- (6,7.25) -- cycle}{0.25}
\fxvdr{(3,8) -- (4,8) -- (4,8.116025) -- (3,8.116025) -- cycle}{(3,8) -- (4,8) -- (4,8.116025) -- (3,8.116025) -- cycle}{}{0.12}
\fxvdr{(4,8) -- (5,8) -- (5,8.316987) -- (4,8.316987) -- cycle}{(4,8) -- (5,8) -- (5,8.316987) -- (4,8.316987) -- cycle}{(4.808017,8) -- (5,8) -- (5,8.080037) -- cycle}{0.32}
\fxvdr{(5,8) -- (6,8) -- (6,8.433013) -- (5,8.433013) -- cycle}{(5,8) -- (6,8) -- (6,8.433013) -- (5,8.433013) -- cycle}{(5,8) -- (6,8) -- (6,8.379838) -- (5,8.119447) -- cycle}{0.43}
\fxvdr{(6,8) -- (7,8) -- (7,8.433013) -- (6,8.433013) -- cycle}{(6,8) -- (7,8) -- (7,8.433013) -- (6,8.433013) -- cycle}{(6,8) -- (7,8) -- (7,8.326195) -- (6,8.380865) -- cycle}{0.43}
\fxvvel{2}{3.5}{7}{-0.128}
\fxvvel{2}{4.5}{7}{-0.350}
\fxvvel{2}{5.5}{7}{-0.477}
\fxvvel{2}{6.5}{7}{-0.477}
\fxvvel{2}{3.5}{8}{-0.222}
\fxvvel{2}{4.5}{8}{-0.605}
\fxvvel{2}{5.5}{8}{-0.827}
\fxvvel{2}{6.5}{8}{-0.827}
\fxvvel{2}{3.5}{9}{-0.256}
\fxvvel{2}{4.5}{9}{-0.699}
\fxvvel{2}{5.5}{9}{-0.955}
\fxvvel{2}{6.5}{9}{-0.955}
\fxvdil{6}{6}{+0.250}
\fxvdil{5}{7}{+0.183}
\fxvdil{5}{6}{+0.250}
\fxvdil{6}{7}{+0.183}
\fxvdil{4}{7}{+0.134}
\fxvdil{4}{6}{+0.183}
\fi
\ifnum\fxvpanel=3
\fxvcell{3}{6}{4}{7}{0.64}{0}{0}
\fxvliquid{(3,6) -- (4,6) -- (4,7) -- (3.857511,7) -- (3,6.170542) -- cycle}
\fxvplic{3.857511}{7}{3}{6.170542}
\fxvcell{4}{6}{5}{7}{1.00}{1}{0}
\fxvliquid{(4,6) -- (5,6) -- (5,7) -- (4,7) -- cycle}
\fxvcell{5}{6}{6}{7}{1.00}{1}{0}
\fxvliquid{(5,6) -- (6,6) -- (6,7) -- (5,7) -- cycle}
\fxvcell{6}{6}{7}{7}{1.00}{1}{0}
\fxvliquid{(6,6) -- (7,6) -- (7,7) -- (6,7) -- cycle}
\fxvcell{3}{7}{4}{8}{0.09}{0}{0}
\fxvliquid{(3.505936,7) -- (4,7) -- (4,7.383694) -- cycle}
\fxvplic{3.505936}{7}{4}{7.383694}
\fxvcell{4}{7}{5}{8}{0.52}{1}{0}
\fxvliquid{(4,7) -- (5,7) -- (5,7.904064) -- (4,7.127455) -- cycle}
\fxvplic{5}{7.904064}{4}{7.127455}
\fxvcell{5}{7}{6}{8}{0.82}{1}{0}
\fxvliquid{(5,7) -- (6,7) -- (6,7.967065) -- (5,7.666194) -- cycle}
\fxvplic{6}{7.967065}{5}{7.666194}
\fxvcell{6}{7}{7}{8}{0.92}{1}{0}
\fxvliquid{(6,7) -- (7,7) -- (7,7.972461) -- (6,7.868574) -- cycle}
\fxvplic{7}{7.972461}{6}{7.868574}
\fxvcell{3}{8}{4}{9}{0.00}{0}{0}
\fxvcell{4}{8}{5}{9}{0.00}{0}{0}
\fxvcell{5}{8}{6}{9}{0.00}{0}{0}
\fxvcell{6}{8}{7}{9}{0.00}{0}{0}
\fxvdil{6}{6}{-0.250}
\fxvdil{5}{7}{-0.183}
\fxvdil{5}{6}{-0.250}
\fxvdil{6}{7}{-0.183}
\fxvdil{4}{7}{-0.134}
\fxvdil{4}{6}{-0.183}
\fi
%
    \end{tikzpicture}}%
  \begin{subfigure}[t]{0.32\textwidth}
    \centering
    \fxvpanelpicture{1}
    \caption{$\condfrac{\liq,\timelevel{n}}{c}$; sweep~1 along $\unitvec{1}$.}
    \label{fig:vof:sweeps:start}
  \end{subfigure}%
  \hfill
  \begin{subfigure}[t]{0.32\textwidth}
    \centering
    \fxvpanelpicture{2}
    \caption{$\condfrac{\liq,\splitlevel{1}}{c}$; sweep~2 along $\unitvec{2}$.}
    \label{fig:vof:sweeps:first}
  \end{subfigure}%
  \hfill
  \begin{subfigure}[t]{0.32\textwidth}
    \centering
    \fxvpanelpicture{3}
    \caption{$\condfrac{\liq,\timelevel{n+1}}{c}=\condfrac{\liq,\splitlevel{2}}{c}$.}
    \label{fig:vof:sweeps:second}
  \end{subfigure}
  \caption{One time step of~\cref{eq:vof:sweep} on the uniform fixture of
    \code{figures/vof.jl} (unit cells, the streamfunction velocity of
    \cref{sec:vof:conservation}, $\dt=\dtwy=\chkVofSweepDt$, order
    $(\sweepaxis{1},\sweepaxis{2})=(1,2)$; window
    $3\le\coordk{1}\le7$, $6\le\coordk{2}\le9$). Numbers: conditional
    fractions; dots: frozen indicator $\indicator{\liq}_c=1$, fixed at the
    start of the step; dashed red: the donating rectangles of the next sweep,
    hatched their liquid; arrows: the velocity of the next sweep; purple
    (b,~c): the dilation volume of the preceding sweep per cell volume,
    $\indicator{\liq}_c\sum_{f}\orient{c}{f}\dt\meas{f}\volumeflux^{\open}_f/\meas{c}$,
    which reaches $\chkVofSweepDilation$ in a sweep but cancels over the two
    sweeps to $\chkVofSweepDilationSum$ because $\divh\velo=0$; in full cells
    it balances the divergence of the one-dimensional liquid flux, so they stay
    full. The PLIC lines are reconstructed after each sweep.}
  \label{fig:vof:sweeps}
\end{figure}
